\subsection{Defensive I/O Patterns}

Defensive I/O code accepts that the filesystem is external state. A path may exist, disappear, become inaccessible, become a directory, or be changed by another process. The program therefore needs patterns that handle failure cleanly instead of assuming that every file operation will succeed.

Python commonly uses two styles: try the operation and handle the exception, or check conditions before attempting the operation. Both are useful, but they do not provide the same guarantees.

\subsubsection{EAFP File Access: Trying the Operation and Handling the Exception}

EAFP means ``easier to ask forgiveness than permission''. In file access, this means the program attempts the operation directly and catches the exception if the operating system refuses it.

\begin{verbatim}
from pathlib import Path

path = Path("missing_eafp_demo.txt")

print("Path object:")
print(f"path       = {path}")
print(f"type(path) = {type(path)}")

print()
print("Selected names from dir(path):")
for name in ["exists", "is_file", "is_dir", "name", "parent", "resolve"]:
    print(f"{name:10} -> {name in dir(path)}")

print()
print("EAFP pattern: try to read, then handle failure.")

try:
    with open(path, "r", encoding="utf-8") as file_obj:
        content = file_obj.read()

    print("File content:")
    print(content)

except FileNotFoundError as err:
    print("The file does not exist.")
    print(f"type(err) = {type(err)}")
    print(f"err       = {err}")

except PermissionError as err:
    print("The file exists, but access was denied.")
    print(f"type(err) = {type(err)}")
    print(f"err       = {err}")

except OSError as err:
    print("Another operating-system I/O error occurred.")
    print(f"type(err) = {type(err)}")
    print(f"err       = {err}")
\end{verbatim}

This pattern is strong because the attempted operation is the same operation whose failure is being handled. The program does not merely ask whether the file looked usable a moment earlier. It asks the operating system to actually perform the read.

EAFP is especially useful when failure can happen for many external reasons. The path may be missing, permission may be denied, the path may point to a directory, the file may be locked, or the filesystem may reject the operation.

A successful EAFP example looks ordinary:

\begin{verbatim}
from pathlib import Path

path = Path("existing_eafp_demo.txt")

with open(path, "w", encoding="utf-8") as file_obj:
    file_obj.write("Serenity now!\n")
    file_obj.write("The answer is 42.\n")

try:
    with open(path, "r", encoding="utf-8") as file_obj:
        content = file_obj.read()

    print("Read succeeded:")
    print(content)

    print("Object inspection:")
    print(f"type(content) = {type(content)}")
    print(f"len(content)  = {len(content)}")

except OSError as err:
    print("Read failed.")
    print(f"type(err) = {type(err)}")
    print(f"err       = {err}")
\end{verbatim}

The \texttt{except OSError} branch is broader than \texttt{FileNotFoundError} or \texttt{PermissionError}. This is useful when the program wants to handle general filesystem failure at one boundary.

\subsubsection{LBYL File Access: Checking Conditions Before Opening}

LBYL means ``look before you leap''. In file access, this means the program checks conditions before attempting the operation.

\begin{verbatim}
from pathlib import Path

path = Path("lbyl_demo.txt")

with open(path, "w", encoding="utf-8") as file_obj:
    file_obj.write("No soup for you!\n")

print("LBYL pattern: check before opening.")

exists_result = path.exists()
file_result = path.is_file()

print(f"path.exists()  = {exists_result}")
print(f"path.is_file() = {file_result}")

print()
print("Object types:")
print(f"type(exists_result) = {type(exists_result)}")
print(f"type(file_result)   = {type(file_result)}")

print()
print("Selected names from dir(exists_result):")
for name in ["bit_length", "to_bytes", "real", "imag"]:
    print(f"{name:12} -> {name in dir(exists_result)}")

if path.exists() and path.is_file():
    with open(path, "r", encoding="utf-8") as file_obj:
        content = file_obj.read()

    print()
    print("File content:")
    print(content)
else:
    print("The path is not currently a readable file candidate.")
\end{verbatim}

The checks are clear and readable. They are useful when the program wants to choose different branches based on current filesystem state.

However, the checks do not open the file. They only report what appeared to be true when the checks were made. The later \texttt{open()} call is still a separate operation that can fail.

LBYL is also useful when producing clearer messages for a human user:

\begin{verbatim}
from pathlib import Path

path = Path("lbyl_missing_demo.txt")

if not path.exists():
    print(f"Cannot read {path}: the path does not exist.")
elif not path.is_file():
    print(f"Cannot read {path}: the path is not a file.")
else:
    with open(path, "r", encoding="utf-8") as file_obj:
        content = file_obj.read()

    print(content)
\end{verbatim}

This style makes the decision tree visible. It is often helpful in scripts, teaching examples, command-line utilities, and validation code.

The weakness is that the checks and the operation are not one indivisible action.

\subsubsection{Atomicity Concerns: Why Existence Checks Can Become Invalid Before Use}

The filesystem can change between two Python statements. Another process can create, delete, rename, or replace a file after the program checks it but before the program opens it. This is often called a check-then-use problem.

The following example shows the shape of the issue without needing another process.

\begin{verbatim}
from pathlib import Path

path = Path("atomicity_demo.txt")

with open(path, "w", encoding="utf-8") as file_obj:
    file_obj.write("This file exists for the first check.\n")

print("First check:")
print(f"path.exists() = {path.exists()}")

print()
print("Simulating an external change by deleting the file before open().")
path.unlink()

print(f"path.exists() = {path.exists()}")

print()
print("Now the earlier check is no longer useful:")

try:
    with open(path, "r", encoding="utf-8") as file_obj:
        content = file_obj.read()

    print(content)

except FileNotFoundError as err:
    print("The file disappeared before it was opened.")
    print(f"type(err) = {type(err)}")
    print(f"err       = {err}")
\end{verbatim}

The first existence check was true, but it did not reserve the file. It did not lock the file. It did not guarantee that a later read would succeed. The external state changed before the actual operation.

This is why EAFP is often preferred for the final file operation. The program still may use LBYL for user-friendly diagnostics, but the real operation should remain protected by exception handling.

A safer creation pattern uses exclusive creation mode \texttt{"x"}. It asks the operating system to create the file only if it does not already exist.

\begin{verbatim}
from pathlib import Path

path = Path("exclusive_create_demo.txt")

print("Trying exclusive creation:")

try:
    with open(path, "x", encoding="utf-8") as file_obj:
        file_obj.write("Created exactly once. D'oh!\n")

    print("File was created.")

except FileExistsError as err:
    print("File already exists; it was not overwritten.")
    print(f"type(err) = {type(err)}")
    print(f"err       = {err}")

print()
print("Trying exclusive creation again:")

try:
    with open(path, "x", encoding="utf-8") as file_obj:
        file_obj.write("This should not replace the old file.\n")

    print("File was created.")

except FileExistsError as err:
    print("Second creation failed safely.")
    print(f"type(err) = {type(err)}")
    print(f"err       = {err}")
\end{verbatim}

This is better than checking \texttt{exists()} and then opening with \texttt{"w"}, because \texttt{"w"} would replace existing content. The mode \texttt{"x"} combines the existence condition with the creation operation.

The practical rule is:

\begin{itemize}
    \item use EAFP when performing the actual file operation;
    \item use LBYL when clear pre-checks help produce better messages;
    \item do not treat existence checks as permanent guarantees;
    \item use operation-level safeguards such as \texttt{"x"} when overwriting would be unsafe.
\end{itemize}

Filesystem state belongs to the operating system and may be shared with other processes. Defensive I/O code must therefore protect the actual operation, not only the condition checked before it.